\subsection{Operations on the product of two measures}

PROPOSITION 10. --- Let \(g\) (resp. \(g'\) ) be a complex function (or a function with values in \(\overline{\mathbf{R}}\) ) defined on \(T\) (resp. \(T'\) ).

\begin{enumerate}
\def\labelenumi{\alph{enumi})}
\tightlist
\item
  If \(g\) (resp. \(g'\) ) is locally integrable for \(\mu\) (resp. \(\mu'\) ), then the function \(g \otimes g': (t, t') \mapsto g(t)g'(t')\) is locally integrable for \(\nu = \mu \otimes \mu'\) , and
\end{enumerate}

\[
(g \cdot \mu) \otimes (g ^ {\prime} \cdot \mu^ {\prime}) = (g \otimes g ^ {\prime}) \cdot (\mu \otimes \mu^ {\prime}).\tag{15}
\]

\begin{enumerate}
\def\labelenumi{\alph{enumi})}
\setcounter{enumi}{1}
\item
  Conversely, if \(g \otimes g'\) is locally \(\nu\) -integrable, and if \(g'\) is not locally \(\mu'\) -negligible, then \(g\) is locally \(\mu\) -integrable.
\item
  Let K and \(K'\) be two compact subsets of T and \(T'\) , respectively; Cor. 2 of Prop. 8 shows that the function \((t, t') \mapsto g(t)g'(t')\varphi_{K \times K'}(t, t')\) , equal to \((g\varphi_{K}) \otimes (g'\varphi_{K'})\) , is \(\nu\) -integrable. Consequently, \(g \otimes g'\) is locally \(\nu\) -integrable. One then verifies immediately that the second member of (15) satisfies the characteristic property of product measures (Ch. III, §4, No.~1, Th. 1).
\item
  Now suppose that \(g \otimes g'\) is locally \(\nu\) -integrable, and that \(g'\) is not locally \(\mu'\) -negligible. Let \(\mu_1\) be a positive measure with compact support such that \(\mu_1 \leqslant \mu\) ; \(g \otimes g'\) being \((\mu_1 \otimes \mu')\) -measurable, \(t \mapsto g(t)g'(t')\) is \(\mu_1\) -measurable except for a locally \(\mu'\) -negligible set of values of \(t'\) (Prop. 2).
\end{enumerate}

Since \(g'\) is not zero locally \(\mu'\) -almost everywhere, from this we deduce that that g is \(\mu_{1}\) -measurable, then \(\mu\) -measurable on decomposing \(\mu\) into a sum of a family of measures with compact support ( \(\S2\) , No.~3, Prop. 4 and \(\S2\) , No.~2, Prop. 2). Having established this point, we may reduce to the case that g and \(g'\) are \(\geqslant0\) , on replacing g and \(g'\) by their absolute values if necessary. Let K be any compact subset of T, and let \(K'\) be a compact subset of \(T'\) such that \(\int g'\varphi_{K'}d\mu'\neq0\) . By Prop. 8,

\[
\left(\int^ {\bullet} g \varphi_ {K} d \mu\right) \left(\int^ {\bullet} g ^ {\prime} \varphi_ {K ^ {\prime}} d \mu^ {\prime}\right) = \iint^ {\bullet} (g \otimes g ^ {\prime}) \varphi_ {K \times K ^ {\prime}} d \mu d \mu^ {\prime} <   + \infty .
\]

The first factor of the first member is therefore finite, and this completes the proof.

This proposition extends to complex measures, thanks to Prop. 3 of Ch. III, §4, No.~2.

PROPOSITION 11. --- Let \(\pi\) (resp. \(\pi'\) ) be a mapping of T (resp. T') into a locally compact space T \(_1\) (resp. T \(_1'\) ).

\begin{enumerate}
\def\labelenumi{\alph{enumi})}
\item
  If \(\pi\) (resp. \(\pi'\) ) is \(\mu\) -proper (resp. \(\mu'\) -proper), then the mapping \(\pi \times \pi'\) is \((\mu \otimes \mu')\) -proper and \((\pi \times \pi')(\mu \otimes \mu') = \pi(\mu) \otimes \pi'(\mu')\) .
\item
  Conversely, if \(\pi \times \pi'\) is \((\mu \otimes \mu')\) -proper and \(\mu' \neq 0\) , then \(\pi\) is \(\mu\) -proper.
\item
  For, \(\pi \times \pi'\) is \((\mu \times \mu')\) -measurable by Cor. 1 of Prop. 3 of No.~2. On the other hand, if K (resp. \(K'\) ) is a compact subset of \(T_1\) (resp. \(T_1'\) ), then \(\overline{\pi}^1(K)\) and \(\overline{\pi}'(K')\) are essentially integrable for \(\mu\) and \(\mu'\) , respectively, therefore \(\overline{\pi}^1(K) \times \overline{\pi}'(K')\) is essentially integrable for \(\mu \otimes \mu'\) (Cor. 2 of Prop. 8). This proves that \(\pi \times \pi'\) is \((\mu \times \mu')\) -proper. Now let \(\mu_1 = \pi(\mu)\) , \(\mu_1' = \pi'(\mu')\) , \(\nu_1 = (\pi \times \pi')(\mu \otimes \mu')\) ; for \(f \in \mathcal{K}(T_1)\) and \(f' \in \mathcal{K}(T_1')\) , we have
\end{enumerate}

\[
\begin{array}{r l} \iint f (\pi (t)) f ^ {\prime} (\pi^ {\prime} (t ^ {\prime})) d \mu (t) d \mu^ {\prime} (t ^ {\prime}) \\ & = \left(\int f (\pi (t)) d \mu (t)\right) \left(\int f ^ {\prime} (\pi^ {\prime} (t ^ {\prime})) d \mu^ {\prime} (t ^ {\prime})\right) \end{array}
\]

(Cor. 2 of Prop. 8), which proves that \(\nu_{1} = \mu_{1}\otimes \mu_{1}^{\prime}\) (Ch. III, §4, No.~1, Th. 1).

\begin{enumerate}
\def\labelenumi{\alph{enumi})}
\setcounter{enumi}{1}
\tightlist
\item
  Now suppose that \(\pi \times \pi'\) is \(\mu \otimes \mu'\) -proper and that \(\mu' \neq 0\) . Let \(\mu_1\) be a measure \(\leqslant \mu\) with compact support. The function \(\pi \times \pi'\) being measurable for \(\mu_1 \otimes \mu'\) , the mapping \(t \mapsto (\pi(t), \pi'(t'))\) is \(\mu\) -measurable except for \(t'\) forming a locally \(\mu'\) -negligible set (No.~2, Prop. 2). Since \(\mu' \neq 0\) , it follows that \(\pi\) is \(\mu_1\) -measurable, and finally that \(\pi\) is \(\mu\) -measurable (§2, No.~3, Prop. 4 and §2, No.~2, Prop. 2). It remains to show that \(\mu^{\bullet}(f \circ \pi) < +\infty\) for every function \(f \in \mathcal{K}_{+}(\mathrm{T}_{1})\) . If \(\mu\) is zero, this property is obvious. If \(\mu\) is not zero, then neither is \(\mu \otimes \mu'\) , consequently \((\pi \times \pi') (\mu \otimes \mu') \neq 0\) (§6, No.~2, Cor. 1 of Prop. 2). By Lemma 1 of Ch. III, §4, No.~1, there exist two functions \(g \in \mathcal{K}_{+}(\mathrm{T}_{1})\) , \(g' \in \mathcal{K}_{+}(\mathrm{T}_{1}')\) such that
\end{enumerate}

\[
\langle (\pi \times \pi^ {\prime}) (\mu \otimes \mu^ {\prime}), g \otimes g ^ {\prime} \rangle \neq 0.
\]

This expression being equal to \(\langle\mu\otimes\mu',(g\circ\pi)\otimes(g'\circ\pi')\rangle\) by the definition of image measures, Prop. 8 implies that \(\mu^{\bullet}(g'\circ\pi')\neq0\) . We then have, by Prop. 8 and by Prop. 2 of §6, No.~2,

\[
\begin{array}{r l} \left(\int^ {\bullet} (f \circ \pi) d \mu\right) & \left(\int^ {\bullet} (g ^ {\prime} \circ \pi^ {\prime}) d \mu^ {\prime}\right) = \iint^ {\bullet} (f \circ \pi) \otimes (g ^ {\prime} \circ \pi^ {\prime}) d \mu d \mu^ {\prime} \\ & = \iint^ {\bullet} (f \otimes g ^ {\prime}) d ((\pi \times \pi^ {\prime}) (\mu \otimes \mu^ {\prime})) <   + \infty . \end{array}
\]

The first integral in the first member is therefore finite, which completes the proof.

This result extends at once to the product of two complex measures (apply the statement to their absolute values). The same is true for the following proposition.

PROPOSITION 12. --- Let X (resp. X') be a locally compact subspace of T (resp. T'). Then, the induced measure \((\mu \otimes \mu')_{X \times X'}\) on the locally compact subspace \(X \times X'\) of \(T \times T'\) is equal to the product \(\mu_{X} \otimes \mu'_{X'}\) of the measures induced on X and \(X'\) by \(\mu\) and \(\mu'\) , respectively.

For, if \(f \in \mathcal{K}(\mathrm{X})\) and \(f' \in \mathcal{K}(\mathrm{X}')\) , then

\[
\iint_ {\mathrm{X} \times \mathrm{X} ^ {\prime}} f (t) f ^ {\prime} (t ^ {\prime}) d \mu (t) d \mu^ {\prime} (t ^ {\prime}) = \left(\int_ {\mathrm{X}} f (t) d \mu (t)\right) \left(\int_ {\mathrm{X} ^ {\prime}} f ^ {\prime} (t ^ {\prime}) d \mu^ {\prime} (t ^ {\prime})\right)
\]

by Cor. 2 of Prop. 8, which proves, by the definition of induced measures (Ch. IV, §5, No.~7) that

\[
(\mu \otimes \mu^ {\prime}) _ {\mathbf {X} \times \mathbf {X} ^ {\prime}} = \mu_ {\mathbf {X}} \otimes \mu_ {\mathbf {X} ^ {\prime}} ^ {\prime}
\]

(Ch. III, §4, No.~1, Th. 1).

